\hypertarget{the-reference-monitor}{%
\chapter{2. The reference monitor}\label{the-reference-monitor}}

Host security begins at the first instruction the processor runs. From
there a chain is built upward: firmware, boot loader, kernel, the login
that authenticates a person, each link measured or signed by the one
before it, decades and fortunes spent on making each hard enough to
stand on. It is not perfect, but it is real, and the effort behind it is
enormous, because rooting trust in hardware and carrying it up through
cryptography is hard. Where that chain reaches us it has settled one
thing we can build on: the operator is the only one who can normally
take on their uid. The host beneath us, and everything outside your own
home, we take as given. We do not rebuild the chain. We inherit it.

Our work begins where the chain ends. It reaches the uid and stops. It
settles who you are and says nothing about the code now running as you
(Chapter 1). The same confidence has to be carried further down, into
the workload itself, and that is the hard part. The host's root of trust
is something we stand on, not something we hold: every privileged thing,
every layer hardened over those decades, belongs to the chain we
inherited, not to us. Beneath it there is no second foundation, only the
home you own and a uid that, by the host's design, may do anything you
may do. What we build has to reach down into the workload from a root we
did not lay, with no privilege of its own.

A reference monitor (\texttt{reference-monitor}) is the structure for
it. Anderson named it in 1972 {[}Anderson{]}, in the report that set out
how a system could enforce a policy it was entitled to trust. It is the
part of a system that every access passes through, permitting or denying
each against the policy, and its value is less the checkpoint than the
three properties Anderson required before anyone relies on one. Those
properties have been the measure of an access-control mechanism ever
since {[}Irvine{]}. It must mediate completely: every access, with no
path around. It must be tamperproof: the code it constrains cannot alter
it or switch it off. And it must be verifiable: small and simple enough
that its correctness can be checked. Miss any one and the other two stop
meaning anything. A complete, tamperproof monitor no one can verify is
trusted on faith; a complete, verifiable monitor that can be switched
off is a suggestion. The count is three in Anderson and still three in
Irvine, whose first property, always invoked, means every access is
mediated, since whatever goes unmediated can be bypassed {[}Irvine{]}. A
later teaching mnemonic, NEAT, splits that property in two and counts
four: non-bypassable, evaluable, always invoked,
tamper-proof.\footnote{See \emph{Reference monitor}, Wikipedia.
  \url{https://en.wikipedia.org/wiki/Reference_monitor}} What reads here
as complete mediation is that first property whole, the monitor on every
path and impossible to route around, kept as one because a path it never
sees and a path that slips around it are the same hole, closed below by
the same means. The model has a formal account, too. Schneider showed in
2000 that a monitor which clears each access as it runs enforces exactly
the safety properties, those that forbid something bad from happening,
with complete mediation and the monitor's isolation from the code it
watches among the conditions it requires {[}Schneider{]}.

Anderson's monitor was itself the root of trust: privileged, at the
foundation, protecting the system from its users. Kennel's root of trust
is not the monitor. It is the operator, handed up to us by the host's
chain, and the monitor sits below it (unprivileged, because the
privilege is all upstream), carrying that trust down into the code the
operator runs. The three properties are unchanged. The trust root has
moved out from under the monitor to above it, and that shift is the
whole of the difficulty.

\hypertarget{complete-mediation}{%
\section{2.1 Complete mediation}\label{complete-mediation}}

Complete mediation {[}SaltzerSchroeder{]} is the property that a
resource cannot be reached except through the monitor. Not most ways in.
Every way. A monitor that holds the front door and misses a side window
has not narrowed the attack, it has moved it, and a workload under
optimisation pressure (Chapter 1) finds the window faster than the
designer found the door.

This is the property the familiar boundaries give up. A sandbox or a
container decides what is reachable and then withdraws; once a thing is
reachable, what is done with it goes unwatched. Complete mediation is
the opposite demand: that use itself is what the monitor governs, every
time, for as long as the workload runs (\texttt{mediate-use-not-reach}).

The difficulty is that completeness is a property of the whole surface
at once, and the surface is wide. A workload reaches the world through
many kinds of resource (files, connections, processes, services,
devices), and each kind is a place a path might be left open. There are
two ways to close a path, and only two. The resource can be removed, so
there is no access to mediate because there is nothing there
(\texttt{construction-by-absence}). Or it can be interposed upon, so
every use is a request the monitor answers
(\texttt{interpose-as-transaction}). The next chapter is those two
limbs, and why between them they cover the surface. What matters here is
the standard they serve: every path is either absent or mediated, with
no third state in which a resource is present, reachable, and
ungoverned.

\hypertarget{tamperproof}{%
\section{2.2 Tamperproof}\label{tamperproof}}

A monitor the workload can switch off is not a monitor. In Chapter 1 an
agent, blocked by its sandbox, reasoned its way to disabling the sandbox
and carrying on (\texttt{T2.1}). It held the controls of its own
confinement, so its confinement was advisory. Tamperproofness is the
requirement that this cannot happen: nothing the workload can do reaches
the monitor.

For a privileged kernel the property comes close to free. The same
boundary that makes the monitor more privileged also keeps it out of
reach. Kennel has no privilege boundary to lean on. The monitor and the
workload run as the same user, and a user can in general change what
runs as them. So tamperproofness cannot rest on the monitor being
stronger than the workload. It rests on the workload having no handle on
the monitor at all.

Two things make that hold. The first is that the boundary exists before
the workload does. The monitor does not start the workload and then race
it to install controls; it builds the confined world first, and only
then places the workload inside it. There is no moment when running code
can intercept its own confinement, because at the moment of confinement
there is no running code. The second is that the monitor's controls are
not objects the workload can see. The configuration that enforces the
policy lives outside the workload's reach: there is no name it can
resolve to find it and no privilege it can borrow to undo it
(\texttt{good-boy}: the monitor does not rely on the workload choosing
not to try). The workload keeps its full ordinary authority over its own
world. That world simply does not contain the levers of its own cage.

The enforcement is the operating system's, not the monitor's own runtime
vigilance. The monitor does not stand over each call checking it by
hand; before the workload starts, it arranges for the system to deny on
its behalf, through the system's own unbypassable primitives. The
monitor authors the policy and builds the world; the kernel makes the
denials stick. This is what lets the monitor stay small, which is the
next property's concern, and it is what makes every attempt to step
outside run into the floor rather than into a part of the monitor that
might be pried loose.

\hypertarget{verifiable}{%
\section{2.3 Verifiable}\label{verifiable}}

The first two properties are claims, and a claim is worth what it can be
checked for. A monitor can be complete and tamperproof in its design and
wrong in its code, and if the code is too large to read, no one will
find out which. Verifiability is the requirement that the monitor is
small and simple enough for its correctness to be established rather
than argued. This is where doing less stops being taste and becomes a
property the other two depend on (\texttt{do-less}).

Small has a measure here, not just a sentiment. The trusted part, the
code the whole boundary rests on, should be small enough for one person
to read in full, and arranged so that reading it is enough. Two
disciplines hold it there. Each component carries at most one source of
risk, so a flaw in any one stays bounded and none combines the hazards
that escapes are built from (\texttt{rule-of-1}). And the monitor stays
out of the data path: it sets a channel up and brokers it, but the bytes
that flow do not pass through it, so the traffic a workload generates
never enlarges the code that has to be trusted
(\texttt{control-not-data-plane}). A monitor that copied every byte
would grow with its workload and never hold still long enough to be
read.

The point of all of it is to take faith out of the chain. The other
properties say the monitor cannot be bypassed and cannot be disabled.
Verifiability is what makes those statements about something a person
has read, instead of hopes about something too large to.

\hypertarget{the-same-three-properties-held-against-the-alternatives}{%
\section{2.4 The same three properties, held against the
alternatives}\label{the-same-three-properties-held-against-the-alternatives}}

The three properties are a measuring stick, and the existing tools can
be held against it. None of what follows is a claim that they are bad at
what they do; it is that they were built for a different problem, and
each forfeits on the axis the interactive-user problem needs most. The
comparison is of approach, not of feature lists.

The hand-rolled sandbox (bubblewrap, Firejail, a seccomp-and-Landlock
profile written per app) can be mediating and can be tamperproof. Where
it gives way is at the granularity of mediation: it confines by what the
workload can reach, and once a channel is present its use is ungoverned
(\texttt{mediate-use-not-reach}). The policy is assembled by hand for
each application rather than derived, so verifiability erodes not
because the mechanism is large but because the policy is bespoke and
grows without a discipline to hold it small. It is the right shape for
wrapping one known program; it does not answer the case of arbitrary
code arriving under the user's own identity.

The mandatory-access-control layer (SELinux, AppArmor) is
kernel-enforced and genuinely tamperproof, and mediates more than a
reachability sandbox does. Its forfeit is elsewhere: the policy is
system-global and authored by the distribution, not by the operator per
workload, and it governs how the system's services are labelled rather
than confining a single untrusted thing the user just ran. The trusted
surface is the whole labelling model, which is not a thing one operator
reads in an afternoon, and the problem it solves is the system's
integrity, not the interactive user's authority over code running as
them.

The container (Docker, and the rootless engines) isolates by namespace,
which is real, but it draws a line and steps back rather than mediating
across it: inside the box the workload is typically root, and the
boundary is the box wall, not a monitor on each access. The controlling
daemon is root-equivalent, so the trusted base is large and privileged.
It is an isolation model for shipping and running a service, and it
neither keeps the workload off the user's own authority
(\texttt{split-the-uid}) nor governs use once a socket or mount is
inside the box.

The virtual machine, and the microVM built for this exact fear
(Firecracker, gVisor as its in-between cousin), is the strongest of the
four on the first two properties: a separate kernel or a separate
machine mediates completely and is hard to tamper with from inside. The
forfeit is the shape of the boundary against the shape of the problem. A
machine-per-task has no native view of the user's own home, so the
workload's world is either a copy of the operator's state, stale and
needing reconciliation, or a share of it, which is the thing being
defended against. Tasks in separate machines do not compose laterally;
where this design lets one confined service reach another by name (the
mesh, taken up in Parts II and III), the machine model makes that reach
a heavyweight crossing rather than a brokered call. And a boundary stood
up per task carries a standing cost an operator running many short-lived
tasks pays over and over; how small that cost is here, and how much of
the appeal rests on it, is a measurement the second volume makes and
this one only asserts. The VM answer is correct for a hostile tenant on
shared hardware. It fits the interactive user, working with their own
files under their own identity, badly.

The through-line is that the sandbox governs reach not use, the MAC
layer governs the system not the workload, the container isolates rather
than mediates, and the VM confines by separation from everything
including the user's own state. Each is sound for the problem it was
built for. The interactive-user problem, code running as you that you
did not write, is a different one, and holding the three properties
against it is what the rest of this volume does, one resource class at a
time.

The three properties are one requirement in three parts. Complete
mediation puts every access in front of the monitor. Tamperproofness
keeps the workload's hands off it. Verifiability makes the first two
checkable instead of hoped for, and the design earns each separately:
completeness by closing every path, tamperproofness by building the
boundary before the workload and out of its reach, verifiability by
staying small enough to read. Where the boundary cannot be built, the
workload does not start, because an unconfined workload is the one
result the monitor must never produce (\texttt{refuse-to-start}). That
leaves the first property's two methods still to work out: a resource
made absent, or a resource interposed upon. Those are the two limbs, and
they are the next chapter.
